\section{Hypothesis}
\label{sec:hypothesis}

\subsection{Primary Hypothesis}

Multilingual LLMs exhibit a \enquote{Language-as-Authority} bias: when presented with directly conflicting factual claims in different languages, the model will systematically favor the claim written in English (or the query language), regardless of factual accuracy or prior knowledge.

\subsection{Justification and Literature Support}

This hypothesis is motivated by the predominance of English in pre-training data, which leads models to associate English text with higher confidence. Previous studies show that English can act as a \enquote{semantic attractor} during reasoning \cite{li2025} and that models preferentially cite English documents \cite{ki2025}. This bias is expected to extend beyond citation habits into actual factual trust.

\subsection{Experimental Validation and Falsifiability}

The hypothesis can be tested and falsified using a counterbalanced setup. For example, given an English query: if the model selects Claim A in \texttt{[English: Claim A, Korean: Claim B]}, but switches to Claim B in \texttt{[Korean: Claim A, English: Claim B]}, the hypothesis is supported—the choice is driven by the language of the claim. Conversely, if the model consistently selects Claim A in both cases, the hypothesis is falsified, indicating that the decision is guided by prior knowledge or factual correctness rather than language bias.

\subsection{Experimental Controls}

To prevent confounding factors, the following four controls are applied.

\begin{enumerate}
    \item \textit{Ordering bias} is addressed by shuffling which language appears first in the prompt.
    \item \textit{Translation artifacts} are minimized by using uniform declarative templates rather than unstructured web text.
    \item \textit{Prior knowledge bias} is mitigated by balancing conditions so each claim appears equally in both languages.
    \item \textit{Query language alignment} is accounted for by testing both English and Korean queries.
\end{enumerate}